\section{Experimental Design}
\label{sec:experiments}

\subsection{Shared setup and scope}

The working setup retains a Qwen3-1.7B student and four domains:
mathematics, code, instruction following, and science. Math and IF use
domain-specific Qwen3-1.7B RL teachers. Code and science currently share
one Qwen3-4B teacher as a temporary substitute for separate specialists.
There are thus four routed domains but only three distinct teacher
weight sets. The repeated teacher is a useful context comparison, not
two independent observations of teacher diversity.

The online anchor uses equal domain prompt quotas and weights,
$16$ micro-batches of four responses per update, one optimizer update
per fresh rollout batch, and a $4{,}096$-token reference cap.
The working budget is at most $500$ updates per run, subject to a measured
throughput pilot on the available $96$-GB learner and $48$-GB
rollout/teacher GPUs. Exact revisions, optimizer and decoding settings,
data splits, and evaluation versions remain to be frozen. The repository
does not yet contain measured outcomes for this study.

The core design reuses runs across questions. A representative
single-teacher setting is selected before examining results; it and the
joint setting each compare PG with top-$64$. The joint top-$64$ run with
domain-response averaging also anchors the normalization comparison,
which adds global-token and domain-token averaging. This gives six
training configurations for an initial study, not a requirement to
train a new expert for every teacher pair. Additional single-teacher
settings and repeated seeds broaden the scope only as resources allow.

\subsection{Study 1: length weighting as an explanatory observation}

Use the same top-$64$ loss, teachers, prompt quotas, and initialization
for all three reductions in Section~\ref{sec:normalization}.
First compare them on identical cached batches: record response counts,
lengths, domain token shares, response-mean gradients, and the covariance
correction. Verify the finite-batch decomposition and measure each
proposed update on a common held-out bank. A mean of unequal micro-batch
means need not equal the intended batch mean, so the implementation
records denominators across accumulation and distributed workers.

Then compare the three online training branches and their per-domain
capability. Their prompt schedules are paired, but each generates its own
on-policy responses. Equal prompt quotas separate length weighting from
changes in the desired domain prior. An unequal-quota diagnostic first
sets target weights to prompt shares before separately testing uniform
weights. A $1{,}024$ versus $4{,}096$ cap comparison is a supporting
length observation, not the paper's new algorithm. Report truncation
and answer-completion rates and distinguish exposure from compute.

\subsection{Study 2a: verify sparsity in single-teacher and multi-teacher OPD}

For the representative single-teacher and joint runs, save the initial,
early, intermediate, and final student checkpoints. Measure cumulative
parameter displacement from each run's own common initialization and
report threshold curves and update-energy concentration. Repeat using
FP32 master weights when available alongside exported-checkpoint
measurements. The comparison establishes the observed sparsity regime
for our experiments and checks whether it persists under joint training;
it does not claim to discover OPD sparsity for the first time.

\subsection{Study 2b: measure teacher-conditioned overlap inside MOPD}

At each selected \emph{joint} checkpoint $\theta_t$, freeze the student
and optimizer state. For each routed domain, draw a diagnostic batch
from that checkpoint's student policy and compute its teacher-conditioned
raw gradient and proposed optimizer step on a separate copy. Compare
their supports at matched coordinate density, retaining both same-layer
random baselines and per-layer results. Repeat over diagnostic batches
so a single noisy estimate is not mistaken for a stable subnetwork.

The primary figure is a teacher/domain overlap matrix at several MOPD
checkpoints, accompanied by each update's concentration. Independent
single-teacher endpoint supports can be shown as a reference, but are
not substituted for contributions within the joint student: the online
students follow different trajectories. A joint checkpoint delta is
also not assumed to decompose uniquely into teacher-specific deltas.
Gradient and proposed-step masks are labeled accordingly.

\subsection{Study 2c: relate teacher distance to subnetwork difference}

Score every available teacher on a common held-out student-prefix bank.
For each pair, compare its distribution distance with the corresponding
support difference, such as $1-J_{ij}$, at the same checkpoint. Record
both teachers' distances to the current student and their update sizes.
Teacher distribution distance must be evaluated on common contexts;
comparing distributions conditioned on different prompts does not
isolate a teacher difference.

The main analysis relates these common-context teacher distances to
\emph{routed} teacher-conditioned supports from Study 2b. A smaller
same-prefix update comparison shows whether a relation persists when
context is also matched for the updates. This supporting check helps
interpret teacher and domain effects without expanding the paper into
a full causal study. Available intermediate teacher checkpoints can
supply additional distance variation without new teacher training.
If the teacher pool remains small, show the individual pairs and report
a setting-specific relationship rather than a universal monotone law.
Pairs sharing a teacher or training lineage are not independent samples.

The intended figure plots teacher-pair distance against support
difference, with checkpoints, domains, and student-relative gaps
identified. Larger, smaller, or unchanged support differences are all
valid outcomes. Teacher--student gap is additionally compared with
update concentration and observed transfer, but no fixed functional
relationship is assumed in advance.

\subsection{Study 3: vocabulary-level supervision density}

Use the representative single-teacher and joint settings in the
same two-by-two comparison: sampled rollout-token PG versus corrected
teacher top-$64$. All use domain-response averaging, matching
initialization, optimizer, masks, cap, prompt schedule, and update budget.
Document full-vocabulary probability normalization and any clipping.
The practical loss comparison is the primary experiment; different
objectives are not mislabeled as an isolated causal change in density.

At the same checkpoints, add full-vocabulary reverse-KL gradients and
proposed optimizer steps on a small common prefix bank in both the
single-teacher and joint settings. These local measurements directly
compare full-vocabulary and sampled supervision without requiring a
new full-vocabulary online run. Separately compare cumulative-delta
sparsity, update energy concentration, and, in the joint setting,
teacher-conditioned support overlap for the online PG/top-$k$ runs. Report actual step size and capability alongside
sparsity, and plot both exposure and compute axes. If the effects are
small and consistent with no material difference at the measurement
resolution, report that observation with its uncertainty and scope.
A stable, meaningful difference can motivate the optional controls in
Appendix~\ref{app:optional_loss}; they are not prerequisites for the
core paper and are not run merely to explain an absent effect.

\subsection{Evaluation and uncertainty}

Use MATH-500 greedy pass@1, a frozen LiveCodeBench pass@1 slice,
IFBench strict accuracy, and GPQA-Diamond average@4. Record evaluator,
model, and data revisions and use fixed decoding across compared
checkpoints. Report raw per-domain scores and the worst-domain change
with the macro-average. Independently trained OPD references, where
available, are compared at matched domain exposure; lack of transfer
in a single-teacher setting is not itself a multi-teacher failure.

Measure generated tokens, domain responses, updates, and GPU hours
separately. Teacher scoring and diagnostics have recorded overhead.
Key online comparisons should have repeated paired training seeds;
initial one-seed runs are explicitly exploratory. Prompt bootstrap
measures evaluation uncertainty conditional on the models, not
training-seed variation. A failure to detect a difference is not by
itself proof of exact equality; report effect sizes and the range the
measurements can distinguish. No numeric sparsity advantage or
capability improvement is claimed before measurement.

\missing{Add the normalization/capability figure, single-versus-joint
sparsity curves, teacher-conditioned overlap matrices, teacher-distance
scatter plots, and PG/top-$k$ comparison table after collecting data.}
